% !TeX program = pdflatex
\documentclass[11pt,a4paper]{amsart}

\newcommand{\notelanguage}{finnish}
% Shared preamble for course notes and exercise sets.

\usepackage[T1]{fontenc}
\usepackage{lmodern}
\providecommand{\notelanguage}{english}
\usepackage[finnish,english]{babel}
\AtBeginDocument{\expandafter\selectlanguage\expandafter{\notelanguage}}

\newcommand{\notesfinnishlanguage}{finnish}
\ifx\notelanguage\notesfinnishlanguage
\newcommand{\notestheoremname}{Lause}
\newcommand{\noteslemmaname}{Lemma}
\newcommand{\notespropositionname}{Propositio}
\newcommand{\notescorollaryname}{Seurauslause}
\newcommand{\notesdefinitionname}{Määritelmä}
\newcommand{\notesexamplename}{Esimerkki}
\newcommand{\notesproblemname}{Tehtävä}
\newcommand{\notesexercisename}{Tehtävä}
\newcommand{\notesremarkname}{Huomautus}
\newcommand{\notessolutionname}{Ratkaisu}
\else
\newcommand{\notestheoremname}{Theorem}
\newcommand{\noteslemmaname}{Lemma}
\newcommand{\notespropositionname}{Proposition}
\newcommand{\notescorollaryname}{Corollary}
\newcommand{\notesdefinitionname}{Definition}
\newcommand{\notesexamplename}{Example}
\newcommand{\notesproblemname}{Problem}
\newcommand{\notesexercisename}{Exercise}
\newcommand{\notesremarkname}{Remark}
\newcommand{\notessolutionname}{Solution}
\fi

\usepackage[a4paper,margin=30mm]{geometry}
\usepackage{microtype}
\usepackage{graphicx}
\usepackage{enumitem}

\newcommand{\exerciseheading}[2]{%
  \par\noindent
  {\large\bfseries #1\par}
  \vspace{0.1em}
  \noindent{\large #2\par}
  \vspace{1.5em}
}

\usepackage{amsmath}
\usepackage{amssymb}
\usepackage{amsthm}
\usepackage{mathtools}
\usepackage{aliascnt}

% Linear systems: the argument gives the number of alignment pairs.
\newenvironment{linearsystem}[1]
  {\left\{\begin{alignedat}{#1}}
  {\end{alignedat}\right.}

\usepackage[hidelinks,pdfusetitle]{hyperref}

\numberwithin{equation}{section}
\setlist{itemsep=0.25em,topsep=0.5em}

\theoremstyle{plain}
\newtheorem{theorem}{\notestheoremname}[section]
\newaliascnt{lemma}{theorem}
\newtheorem{lemma}[lemma]{\noteslemmaname}
\aliascntresetthe{lemma}
\newaliascnt{proposition}{theorem}
\newtheorem{proposition}[proposition]{\notespropositionname}
\aliascntresetthe{proposition}
\newaliascnt{corollary}{theorem}
\newtheorem{corollary}[corollary]{\notescorollaryname}
\aliascntresetthe{corollary}

\theoremstyle{definition}
\newaliascnt{definition}{theorem}
\newtheorem{definition}[definition]{\notesdefinitionname}
\aliascntresetthe{definition}
\newaliascnt{example}{theorem}
\newtheorem{example}[example]{\notesexamplename}
\aliascntresetthe{example}
\newaliascnt{problem}{theorem}
\newtheorem{problem}[problem]{\notesproblemname}
\aliascntresetthe{problem}
\newtheorem{exercise}{\notesexercisename}

\theoremstyle{remark}
\newaliascnt{remark}{theorem}
\newtheorem{remark}[remark]{\notesremarkname}
\aliascntresetthe{remark}

\newenvironment{solution}
{
\begin{proof}[\notessolutionname]}
  {
\end{proof}}

\usepackage[nameinlink,noabbrev]{cleveref}

\crefname{theorem}{theorem}{theorems}
\crefname{lemma}{lemma}{lemmas}
\crefname{proposition}{proposition}{propositions}
\crefname{corollary}{corollary}{corollaries}
\crefname{definition}{definition}{definitions}
\crefname{example}{example}{examples}
\crefname{problem}{problem}{problems}
\crefname{exercise}{exercise}{exercises}
\crefname{remark}{remark}{remarks}

\ifx\notelanguage\notesfinnishlanguage
\crefname{theorem}{lause}{lauseet}
\crefname{lemma}{lemma}{lemmat}
\crefname{proposition}{propositio}{propositiot}
\crefname{corollary}{seurauslause}{seurauslauseet}
\crefname{definition}{määritelmä}{määritelmät}
\crefname{example}{esimerkki}{esimerkit}
\crefname{problem}{tehtävä}{tehtävät}
\crefname{exercise}{tehtävä}{tehtävät}
\crefname{remark}{huomautus}{huomautukset}
\fi

\newcommand{\NN}{\mathbb{N}}
\newcommand{\NNone}{\mathbb{N}_1}
\newcommand{\ZZ}{\mathbb{Z}}
\newcommand{\QQ}{\mathbb{Q}}
\newcommand{\RR}{\mathbb{R}}
\newcommand{\CC}{\mathbb{C}}
\newcommand{\PP}{\mathbb{P}}

\DeclareMathOperator{\im}{im}
\DeclareMathOperator{\rank}{rank}
\DeclareMathOperator{\syt}{syt}
\DeclarePairedDelimiter{\abs}{\lvert}{\rvert}
\DeclarePairedDelimiter{\norm}{\lVert}{\rVert}
\newcommand{\set}[1]{\left\lbrace#1\right\rbrace}
\newcommand{\maxset}[1]{\max\set{#1}}
\newcommand{\minset}[1]{\min\set{#1}}
\newcommand{\powerset}[1]{\mathcal{P}\!\left(#1\right)}

\usepackage{pgfplots}
\pgfplotsset{compat=1.18}

\title{Lineaarialgebra ja matriisilaskenta I\\Harjoituksen 1 ratkaisut}
\author{}
\date{}

\begin{document}

\exerciseheading{Lineaarialgebra ja matriisilaskenta I}{Harjoituksen
1 ratkaisut}

\begin{exercise}
  Tutki, onko yhtälöryhmällä
  \[
    \text{(a)}\;
    \begin{linearsystem}{2}
      2x_1 &{}-{}& x_2 &= 3 \\
      -4x_1 &{}+{}& 2x_2 &= -6
    \end{linearsystem}
    \quad
    \text{(b)}\;
    \begin{linearsystem}{2}
      -x_1 &{}+{}& 3x_2 &= -2 \\
      2x_1 &{}+{}& x_2 &= -3
    \end{linearsystem}
    \quad
    \text{(c)}\;
    \begin{linearsystem}{2}
      -2x_1 &{}+{}& 6x_2 &= -1 \\
      x_1 &{}-{}& 3x_2 &= 1
    \end{linearsystem}
  \]
  tasan yksi, ei yhtään vai ääretön määrä ratkaisuja. Tulkitse jokaisen
  yhtälöparin ratkaiseminen graafisesti \(x_1x_2\)-tasossa.
\end{exercise}

\begin{solution}
  \leavevmode\par
\begin{enumerate}[label=\alph*), nosep]
  \item Kertomalla yhtälöryhmän toinen yhtälö luvulla \(-\frac12\) saadaan
    \[
      \begin{linearsystem}{2}
        2x_1 &{}-{}& x_2 &= 3 \\
        2x_1 &{}-{}& x_2 &= 3.
      \end{linearsystem}
    \]
    Molemmat yhtälöt kuvaavat siis samaa suoraa \(x_2 = 2x_1 - 3\).
    Suorat ovat päällekkäiset, ja jokainen tämän suoran piste on
    yhtälöryhmän ratkaisu. Ratkaisuja on siten äärettömän monta.

    Valitaan \(x_1=t\in\RR\). Tällöin \(x_2=2t-3\), joten
    \[
      \begin{pmatrix} x_1 \\ x_2
      \end{pmatrix}
      =
      \begin{pmatrix} t \\ 2t-3
      \end{pmatrix}
      =
      \begin{pmatrix} 0 \\ -3
      \end{pmatrix}
      + t
      \begin{pmatrix} 1 \\ 2
      \end{pmatrix}.
    \]
    Ratkaisujoukko on siis
    \[
      \left\{
        \begin{pmatrix} 0 \\ -3
        \end{pmatrix}
        + t
        \begin{pmatrix} 1 \\ 2
        \end{pmatrix}
        \colon\ t \in \RR
      \right\}.
    \]
    \begin{center}
      \begin{tikzpicture}
        \begin{axis}[
            width=7cm, height=6cm,
            axis lines=middle, axis equal image,
            axis line style={thin, gray},
            xmin=-1, xmax=4, ymin=-4, ymax=4,
            xlabel={$x_1$}, ylabel={$x_2$},
            xtick={1.5}, xticklabels={$\frac32$}, ytick={-3},
            tick label style={font=\small},
            domain=-0.5:3.5, samples=2,
            clip=false,
          ]
          \addplot[black, semithick] {2*x-3};
          \node[anchor=west, font=\small] at (axis cs:2.7,2.4)
          {$x_2=2x_1-3$};
        \end{axis}
      \end{tikzpicture}
    \end{center}
  \item Lisäämällä toiseen yhtälöön ensimmäinen yhtälö kerrottuna
    kahdella saadaan
    \[
      \begin{linearsystem}{2}
        -x_1 &{}+{}& 3x_2 &= -2 \\
        &     & 7x_2 &= -7.
      \end{linearsystem}
    \]
    Siis \(x_2=-1\). Sijoittamalla tämä ensimmäiseen yhtälöön saadaan
    \(-x_1-3=-2\), joten \(x_1=-1\).
    Yhtälöryhmällä on siten täsmälleen yksi ratkaisu, ja ratkaisujoukko on
    \[
      \left\{
        \begin{pmatrix} -1 \\ -1
      \end{pmatrix}\right\}.
    \]
    \begin{samepage}
      Graafisesti yhtälöt kuvaavat suoria
      \[
        x_2=\frac13x_1-\frac23
        \qquad\text{ja}\qquad
        x_2=-2x_1-3.
      \]
      Suorien kulmakertoimet ovat erisuuret, joten ne leikkaavat toisensa
      täsmälleen yhdessä pisteessä \((-1,-1)\).
      \begin{center}
        \begin{tikzpicture}
          \begin{axis}[
              width=8cm, height=6cm,
              axis lines=middle, axis equal image,
              axis line style={thin, gray},
              xmin=-3, xmax=2, ymin=-4, ymax=2,
              xlabel={$x_1$}, ylabel={$x_2$},
              xtick=\empty, ytick=\empty,
              samples=2, clip=false,
            ]
            \addplot[black, semithick, domain=-3:2] {x/3-2/3};
            \addplot[black, semithick, domain=-2.5:0.5] {-2*x-3};
            \addplot[only marks, mark=*, mark size=1.5pt, black]
            coordinates {(-1,-1)};
            \node[anchor=south east, font=\small] at (axis cs:-1,-1)
            {$(-1,-1)$};
            \node[anchor=north west, font=\small] at (axis cs:1,-0.3333)
            {$x_2=\frac13x_1-\frac23$};
            \node[anchor=east, font=\small] at (axis cs:-2.2,1.4)
            {$x_2=-2x_1-3$};
          \end{axis}
        \end{tikzpicture}
      \end{center}
    \end{samepage}
  \item Lisäämällä ensimmäiseen yhtälöön toinen yhtälö kerrottuna
    kahdella saadaan
    \[
      (-2x_1+6x_2)+2(x_1-3x_2)=-1+2,
    \]
    eli \(0=1\), mikä on ristiriita. Yhtälöryhmällä ei siis ole ratkaisuja,
    ja ratkaisujoukko on \(\varnothing\).

    \begin{samepage}
      Graafisesti yhtälöt kuvaavat suoria
      \[
        x_2=\frac13x_1-\frac16
        \qquad\text{ja}\qquad
        x_2=\frac13x_1-\frac13.
      \]
      Suorilla on sama kulmakerroin mutta eri vakiotermit. Ne ovat siis
      erilliset yhdensuuntaiset suorat, eikä niillä ole yhteisiä pisteitä.
      \begin{center}
        \begin{tikzpicture}
          \begin{axis}[
              width=8cm, height=4.5cm,
              axis lines=middle, axis equal image,
              axis line style={thin, gray},
              xmin=-2, xmax=2, ymin=-1.2, ymax=1,
              xlabel={$x_1$}, ylabel={$x_2$},
              xtick=\empty, ytick=\empty,
              domain=-2:2, samples=2, clip=false,
            ]
            \addplot[black, semithick] {x/3-1/6};
            \addplot[black, semithick] {x/3-1/3};
            \node[anchor=west, font=\small] at (axis cs:2,0.5)
            {$x_2=\frac13x_1-\frac16$};
            \node[anchor=north west, font=\small] at (axis cs:-2,-1)
            {$x_2=\frac13x_1-\frac13$};
          \end{axis}
        \end{tikzpicture}
      \end{center}
    \end{samepage}
\end{enumerate}
\end{solution}

\begin{exercise}
Kemiallisia reaktioita kuvaavat yhtälöt pitää tasapainottaa niin, että
niissä on kummallakin puolella sama määrä kunkin alkuaineen atomeja.
Määritä seuraavien reaktioyhtälöiden kertoimet \(x_1\), \(x_2\), \(x_3\)
ja \(x_4\), joiksi etsitään mahdollisimman pienet
positiiviset kokonaisluvut:
\begin{enumerate}[label=(\alph*)]
  \item \(x_1\mathrm{NO}_2 + x_2\mathrm{H}_2\mathrm{O}
    \to x_3\mathrm{HNO}_2 + x_4\mathrm{HNO}_3\),
  \item \(x_1\mathrm{C}_6\mathrm{H}_6 + x_2\mathrm{O}_2
    \to x_3\mathrm{C} + x_4\mathrm{H}_2\mathrm{O}\).
\end{enumerate}
\end{exercise}

\begin{exercise}
Etsi kertoimet \(a\), \(b\) ja \(c\) siten, että funktion
\(f(x) = ax^2 + bx + c\) kuvaaja kulkee pisteiden \((1,2)\),
\((-1,6)\) ja \((2,3)\) kautta.
\end{exercise}

\begin{exercise}
Ratkaise yhtälöryhmä
\[
  \begin{linearsystem}{4}
    x_1 &{}+{}& x_2 &{}+{}& x_3 &{}+{}& x_4 &= 0 \\
    2x_1 &{}+{}& 3x_2 &{}-{}& x_3 &{}-{}& x_4 &= 2 \\
    3x_1 &{}+{}& 2x_2 &{}+{}& x_3 &{}+{}& x_4 &= 5 \\
    3x_1 &{}+{}& 6x_2 &{}-{}& x_3 &{}-{}& x_4 &= 4.
  \end{linearsystem}
\]
\end{exercise}

\begin{exercise}
Ratkaise yhtälöryhmä
\[
  \begin{linearsystem}{6}
    x_1 &{}+{}& 3x_2 &{}-{}& 2x_3 & & &{}+{}&
    2x_5 & & &= 0 \\
    2x_1 &{}+{}& 6x_2 &{}-{}& 5x_3 &{}-{}& 2x_4 &{}+{}&
    4x_5 &{}-{}& 3x_6 &= -1 \\
    & & & & 5x_3 &{}+{}& 10x_4 & &
    &{}+{}& 15x_6 &= 5 \\
    2x_1 &{}+{}& 6x_2 & & &{}+{}& 8x_4 &{}+{}&
    4x_5 &{}+{}& 18x_6 &= 6.
  \end{linearsystem}
\]
\end{exercise}

\begin{exercise}
Olkoon \(n \in \ZZ_+\). Todista monisteen Lemma 1.3.2:
yhtälöryhmillä
\[
  \begin{linearsystem}{1}
    a_{11}x_1 + \dots + a_{1n}x_n &= b_1 \\
    a_{21}x_1 + \dots + a_{2n}x_n &= b_2
  \end{linearsystem}
\]
ja
\[
  \begin{linearsystem}{1}
    a_{11}x_1 + \dots + a_{1n}x_n &= b_1 \\
    (a_{21} + a_{11})x_1 + \dots + (a_{2n} +
    a_{1n})x_n &= b_2 + b_1
  \end{linearsystem}
\]
on sama ratkaisujoukko.
\end{exercise}

\end{document}
